\documentclass{ltc-guji}
\setmainfont{SourceHanSerifSC-Regular.otf}
\关闭分页

\begin{document}
\begin{正文}

% 基本用法：默认圈点
\decorate{天地}默认圈点

% 自定义装饰字符
\decorate[char=、]{宇宙}顿号装饰

% 颜色测试
\decorate[char=。, color=blue]{日月}蓝色圈点

\decorate[char=●, color=green, scale=0.4]{辰宿}绿色实心点

% 位置偏移测试
\decorate[char=。, xshift=-0.3em]{寒来}右偏移

\decorate[char=。, yshift=0.3em]{秋收}下偏移

% 缩放测试
\decorate[char=。, scale=0.5]{闰余}缩小圈点

\decorate[char=●, scale=0.6, color=red]{律吕}红色实心点

% 混合使用
云腾\decorate[char=。]{致雨}露结\decorate[char=、]{为霜}混合装饰

% 改字功能
天\改{地}改字测试

% EmphasisMark 着重号功能
\EmphasisMark{金生}着重号默认

\圈点[color=blue]{丽水}蓝色圈点

\着重号[color=green]{玉出}绿色着重号

% ============================================================================
% 专名号/书名号 (LineMark) - PDF 画线实现
% ============================================================================

% 专名号（直线）- 基本用法
\专名号{司馬遷}默认专名号

% 专名号 - 颜色
\专名号[color=blue]{班固}蓝色专名号

% 书名号（波浪线）- 基本用法
\书名号{史記}默认书名号

% 书名号 - 颜色
\书名号[color=blue]{漢書}蓝色书名号

% 连续词语断开测试：两个书名号之间应有小间距
\书名号{史記}\书名号{漢書}连续书名

% 连续词语断开测试：两个专名号之间
\专名号{司馬遷}\专名号{班固}连续专名

% 混合测试：专名号和书名号相邻
\专名号{司馬遷}撰\书名号{史記}混合

% 波浪线振幅测试（标准样式 - 默认）
\书名号[amplitude=small]{小振幅}小

\书名号[amplitude=medium]{中振幅}中

\书名号[amplitude=large]{大振幅}大

% 波浪线样式测试：草书风格
\书名号[style=cursive]{草书波浪}草书

\书名号[style=cursive, amplitude=large]{大草书}大草

% 英文别名兼容性
\Underline{下划线}兼容

\WavyUnderline{波浪线}兼容

% ============================================================================
% 特殊环境中的专名号/书名号
% ============================================================================

% 夹注中的专名号（默认 outward 对齐）
前文\夹注{\专名号{司馬遷}著\书名号{史記}}夹注内

% 夹注中的书名号（默认 outward 对齐）
前文\夹注{\书名号{漢書}由\专名号{班固}撰}夹注内

% 夹注 inward 对齐
前文\夹注[align=inward]{\专名号{司馬遷}著\书名号{史記}}向内

% 夹注 center 对齐
前文\夹注[align=center]{\专名号{司馬遷}著\书名号{史記}}居中

% 夹注 left 对齐
前文\夹注[align=left]{\专名号{司馬遷}著\书名号{史記}}居左

% 侧批中的专名号
\侧批{\专名号{杜甫}}诗圣杜甫

% 文本框中的专名号
\文本框{\专名号{李白}\书名号{蜀道難}}

% ============================================================================
% 全局设置命令 (fix #86)
% ============================================================================

% decorateSetup 全局装饰设置
\装饰设置{color=red}
\decorate{全局红}全局红色圈点
\装饰设置{color=black}

% emphasisMarkSetup 全局着重号设置
\着重号设置{color=blue}
\着重号{全局蓝}全局蓝色着重号
\着重号设置{color=black}

% lineMarkSetup 全局线标设置
\线标设置{color=red}
\专名号{全局红}全局红色专名号
\书名号{全局红}全局红色书名号
\线标设置{color=black}

% ============================================================================
% 旁注 (SideText) - 文字左右装饰字 (fix #94)
% 新 API: \旁注{主字}{右内容}{左内容?}
% ============================================================================

% 基本用法：左右都有旁注
\SideText{陽}{配氏}{字辈}左右旁注

% 只有右边
\旁注{甲}{右注}只有右

% 左右都有
\旁注{乙}{右注}{左注}左右旁注

% 长文本旁注
\旁注{戊}{四五六}{一二三}长旁注

% 文本框旁注 - 右边用文本框
\旁注{己}{\文本框[font-size=10pt, grid-height=10pt, height=3]{配氏}}文本框右

% 文本框旁注 - 左右都用文本框
\旁注{庚}{\文本框[font-size=10pt, grid-height=10pt, height=2]{右注}}{\文本框[font-size=10pt, grid-height=10pt, height=2]{左注}}文本框左右

% ============================================================================
% 发圈法 (FaSheng) - 平上去入声调圈点 (fix #92)
% 平=左下 上=左上 去=右上 入=右下
% ============================================================================

% 基本四声 (单字裸标)
\发圈[平]{天}平声
\发圈[上]{地}上声
\发圈[去]{雨}去声
\发圈[入]{雪}入声

% 显式声键
\发圈[声=平]{东}声=平
\发圈[声=上]{南}声=上

% 多字符 (整串套圈)
\发圈[去]{江山}去声多字

% 中文别名 — 圈发
\圈发[平]{春}圈发别名

% 四声中文长名
\发圈[平声]{风}平声
\发圈[上声]{雨}上声
\发圈[去声]{霜}去声
\发圈[入声]{雪}入声

% 颜色调整
\发圈[去, 颜色=red]{字}红去声

% 自定义字符 (实心点)
\发圈[平, 字符=●, 缩放=0.25]{月}实心平

% 调整偏移 (远离字角)
\发圈[去, 偏移=0.6em]{星}远偏移

% 全局设置
\发圈设置{颜色=blue}
\发圈[平]{东}全局蓝
\发圈[去]{风}全局蓝
\发圈设置{颜色=black}

% 批量平行 — 一对一映射
\发圈批{平上去入}{春夏秋冬}四季

% 句中混合 — 与正文交错
天\发圈[平]{涯}远海\发圈[上]{角}行混合

% 单独命令 \平声 \上声 \去声 \入声
\平声{春}\上声{夏}\去声{秋}\入声{冬}单令四声

% 单独命令 + 选项 (颜色、偏移)
\平声[颜色=red]{东}红平
\去声[偏移=0.6em]{南}远去

% 每个调用单独偏移大小 — 不互相影响
\发圈[平, 偏移=0.3em]{近}近偏
\发圈[平, 偏移=0.45em]{中}中偏
\发圈[平, 偏移=0.6em]{远}远偏
\发圈[平]{常}默认 % 应回到 0.45em 默认

% ============================================================================
% 夹注中的发圈 — 验证子栏 (sub_col) 内的圈点位置
% ============================================================================

% 夹注内单声
前\夹注{\发圈[平]{甲}注\发圈[去]{乙}}夹

% 夹注内四声
前\夹注{\平声{春}\上声{夏}\去声{秋}\入声{冬}}夹

% 夹注内多字+发圈
前\夹注{\发圈[平]{江山}秀\发圈[去]{社稷}稳}夹

\end{正文}
\end{document}
